\section{Part A: Review Paper}
\subsection{Abstract}
Traffic engineering (TE) assigns traffic to network resources to meet performance objectives under changing demand. This review compares distributed shortest-path routing and weight optimization with centralized WAN and data-center approaches. It focuses on what each controller observes, what it changes, and how its evaluation treats operational cost and route stability. These distinctions motivate a bounded SmartTE experiment: whether a periodic port-counter controller with a switching margin and minimum dwell changes path stability and application outcomes in a small emulated topology. The experiment is not a proposal for a new class of TE algorithm and its results are not extrapolated to production networks.

\subsection{Keywords}
Traffic engineering; software-defined networking; OpenFlow; multipath routing; route stability; Mininet.

\subsection{Introduction and motivation}
Traffic engineering steers traffic to use available capacity while satisfying service or operational goals. The problem is not simply to find a route: a network may have several feasible paths, but their usefulness depends on current load, traffic demand, measurement freshness, and the cost of changing forwarding state. A design that reacts slowly may leave a bottleneck in place; one that reacts to every noisy measurement may add control traffic or repeatedly move flows.

Data-center fabrics make this tension concrete. Multi-rooted topologies provide multiple paths between hosts, but a path-selection policy determines whether that capacity is used effectively \cite{alfares2008,greenberg2009}. Software-defined networking (SDN) separates control and forwarding responsibilities, while OpenFlow defines a controller-to-switch interface for flow-table operations \cite{rfc7426,openflow13,openflow2008}. Central visibility can coordinate decisions, but it is not free: telemetry, controller processing, rule installation, and transition behavior all matter. The research question for SmartTE is correspondingly narrow: in one two-path emulation, does adding hysteresis and a dwell timer to a counter-based policy reduce successful route changes while retaining useful throughput and delay? The case study measures that implementation under a fixed workload; it does not establish a general result for SDN or TE.

\subsection{Distributed routing and weight-based engineering}
OSPF is a link-state protocol: routers flood topology information and independently calculate shortest paths using configured link weights \cite{rfc2328}. This distributes route computation and failure response, but an operator still has to choose weights that produce desirable network-wide traffic patterns. Fortz and Thorup formulate weight selection as an optimization problem against a demand matrix, illustrating that the configuration problem can be global even when packet forwarding remains distributed \cite{fortz2000}. Their later discussion of traditional IP routing describes the practical role of weight-based TE and its constraints \cite{fortz2002}. Thus ``centralized optimization'' and ``centralized per-flow forwarding'' are different control arrangements: an external optimizer may update costs while each router continues to use its IGP.

Equal-cost multipath (ECMP) uses multiple next hops where shortest-path costs tie. RFC 2992 analyzes hash-threshold selection and the disruption of existing flow assignments when the next-hop set changes \cite{rfc2992}. Hashing scales without a controller decision for each flow, but flow-size imbalance and collisions can make utilization uneven. Directly moving selected flows in response to congestion provides more control, at the cost of requiring adequate measurements and update logic. ECMP's disruption analysis also highlights that a change in route assignments is itself an outcome to measure, not merely an implementation detail.

\subsection{Topology diversity and data-center multipath}
Fat-tree and related multi-rooted designs use commodity switches to create path diversity between hosts \cite{alfares2008}. VL2 develops a scalable data-center architecture and uses Valiant Load Balancing to spread traffic across paths \cite{greenberg2009}. These works establish architectural opportunities rather than a guarantee that arbitrary traffic will distribute evenly. The assignment rule, traffic matrix, flow sizes, and available path information determine whether nominal capacity can be exploited.

Hedera addresses collisions caused by static flow hashing in multi-rooted data-center networks. It detects large flows and centrally schedules their placement across the fabric; its evaluation includes a commodity-switch implementation and a much larger simulated topology than SmartTE \cite{hedera}. This approach uses flow-level information and global placement to avoid overlap, while a simple port-counter controller observes aggregate egress activity and may not know which individual flow caused it. The two systems therefore optimize with different information and at different granularities. Hedera's reported simulation result should not be treated as an expected performance level for a four-switch emulation.

MicroTE studies fine-grained TE using short-term traffic patterns and traffic-matrix information \cite{microte}. Prediction or repeated traffic structure can enable decisions before congestion is fully visible in a counter interval. That advantage depends on the stability and availability of the workload signal; it is not interchangeable with reactive telemetry. For SmartTE, periodic port counters are deliberately simpler and easier to reproduce, but they can be stale by up to a polling interval and include the controlled flow's own traffic. Moving that flow changes the measurement being optimized, creating a feedback loop absent from a static demand estimate.

CONGA offers a useful contrast within congestion-aware data-center routing: it distributes load-balancing decisions in the network rather than relying on a single centralized controller to select a path for each flow \cite{conga}. The approach shows that congestion feedback can be placed at a different point in the system architecture. SmartTE instead uses OpenFlow port statistics collected by OS-Ken and applies one controller policy to the primary flow. These mechanisms are not interchangeable: their measurement locations, switch support, control scope, and evaluated workloads differ.

\subsection{Centralized traffic engineering in WANs}
B4 reports Google's software-defined WAN design and production experience. Its control system combines traffic engineering with application priorities, demand information, and rate control at the network edge \cite{b4}. SWAN also centrally coordinates inter-datacenter traffic, explicitly managing service rates and network configurations; its evaluation combines a testbed prototype with data-driven simulations and studies how to apply changes without transient congestion \cite{swan}. These systems show that centralized TE can coordinate more than route choice: rates, priorities, topology constraints, and update sequencing are part of the policy.

The WAN results are valuable architectural evidence, but their operating assumptions differ from this case study. A provider controls endpoints and traffic classes, has a network-wide demand model, and operates across many links and sites. SmartTE uses one primary TCP flow, one optional UDP competitor, two paths, and periodically sampled counters. It does not allocate WAN service rates, optimize a large demand matrix, or provide a distributed high-availability control plane. Accordingly, B4 and SWAN motivate the value of coordinated state while also showing why results from their production, testbed, and data-driven evaluations cannot be compared numerically with SmartTE.

\subsection{Controller scalability, monitoring, and consistency}
Centralization creates an information advantage but also a scaling question. DevoFlow examines the control-plane and implementation overheads associated with fine-grained flow management and proposes reducing controller involvement for ordinary flows while retaining visibility into significant traffic \cite{devoflow}. This is a different point in the design space from polling a few switch ports once per second, yet the underlying lesson applies: telemetry frequency, message volume, controller work, and update latency are system costs that should be measured. Increasing the polling rate can make the estimate fresher while increasing requests and sensitivity to short-lived variation.

Consistency during a rule update is another distinct concern. Reitblatt and coauthors formalize per-packet and per-flow consistency properties and mechanisms for changing network configurations without exposing packets to an inconsistent mixture of rules \cite{reitblatt}. Those guarantees address how a chosen update is installed. They do not determine how often a TE algorithm chooses a new route. Conversely, a hysteresis margin or dwell timer may reduce decision frequency without proving that a multi-switch update is transiently consistent. SmartTE logs decisions and successful transitions, but it does not claim the formal guarantees of consistent-update mechanisms.

\subsection{Route stability and feedback}
Route stability can refer to several measurable properties: the number of flows remapped after a topology event, repeated oscillation between paths, or transient inconsistency while rules are changed. ECMP disruption analysis treats path remapping after next-hop changes \cite{rfc2992}; hybrid-SDN research studies instability caused by interaction between centralized and distributed control loops \cite{tseng2019}. The latter is not SmartTE's architecture, but it reinforces the need to distinguish a policy's decision dynamics from the correctness of its forwarding updates.

A switching margin requires an alternative path to be better by a minimum utilization difference; a dwell timer prevents another switch until enough time has elapsed. Both can suppress reactions to small or short-lived differences, but both can also delay a useful route change. In a controller whose score includes the primary flow, changing paths modifies the scores: utilization may fall on the old route and rise on the new one. This endogenous feedback can make a plain best-path policy alternate even if exogenous demand is steady. Therefore, successful route changes should be interpreted alongside throughput, round-trip time, loss, response time, and the measurement method. Fewer changes alone do not prove a better service outcome.

\subsection{Evaluation method and emulation limits}
Mininet runs real networking software and applications in a lightweight virtual topology, making it useful for reproducible prototyping \cite{mininet2010}. Mininet-HiFi evaluates fidelity techniques for container-based emulation and discusses the limits of treating emulated performance as equivalent to a physical network \cite{handigol2012}. SmartTE uses Mininet, Open vSwitch, OpenFlow 1.3, and OS-Ken to test controller behavior in a controlled four-switch setup. This supports repeatable comparisons of the declared topology and workload; it does not recreate hardware forwarding, large flow populations, production WAN delay, or controller scaling at fleet size.

Measurement choices also constrain inference. Port-byte deltas estimate link utilization over an interval; iperf3 receiver data measures application throughput; ping supplies sampled delay and loss. These metrics are related but not interchangeable. TCP adjusts its sending behavior in response to congestion, whereas a configured UDP rate offers a different load pattern. Queue limits and host scheduling affect observed outcomes, and short runs can include startup transients. A defensible report therefore names each signal, its clock and aggregation, the number of trials, missing outcomes, and whether intervals are descriptive or part of a statistical test. The M5 report uses five trials per configuration and descriptive 95\% Student-t intervals; it does not claim a hypothesis test.

\subsection{Comparative discussion and research gap}
Table~\ref{tab:related-work} summarizes the contrasts relevant to SmartTE. The methods are complementary rather than direct competitors: they differ in what they observe, what they control, and how their evidence is produced.

\begin{table}[!htbp]\centering\small
\caption{Selected traffic-engineering approaches and their relation to SmartTE.}
\label{tab:related-work}
\begin{tabular}{>{\raggedright\arraybackslash}p{2.35cm}>{\raggedright\arraybackslash}p{3.65cm}>{\raggedright\arraybackslash}p{3.25cm}>{\raggedright\arraybackslash}p{4.6cm}}
\toprule
Approach & Main signal and action & Evaluation focus & Difference from SmartTE \\
\midrule
OSPF weight optimization \cite{fortz2000,fortz2002} & Optimize IGP link weights from a demand matrix; routers then compute shortest paths & Network-wide load under the supplied demand & Changes costs, not per-flow routes \\
ECMP \cite{rfc2992} & Hash each flow over equal-cost next hops & Scaling and path-remapping disruption & No central congestion feedback per flow \\
Hedera / MicroTE / CONGA \cite{hedera,microte,conga} & Schedule large flows, use traffic patterns, or distribute congestion feedback & Data-center capacity under different workloads & More flow/demand context or in-network coordination than SmartTE \\
B4 / SWAN \cite{b4,swan} & Coordinate WAN demand, priorities, rates, and paths & Production or data-driven WAN utilization & Rate control and operational scale are outside SmartTE \\
SmartTE B/C & Periodic port-byte deltas; move one h1--h2 flow & Route changes, TCP, ping, and counters in Mininet & A bounded emulation with aggregate measurements \\
\bottomrule
\end{tabular}
\end{table}

Across these studies, a common theme is that path diversity becomes useful only when traffic assignment matches demand, while the appropriate assignment depends on the available information and control scope. There is no single shared objective: IGP optimization targets network-wide cost, ECMP spreads flows, Hedera schedules large flows, MicroTE exploits traffic patterns, and WAN systems coordinate rates and priorities. Their differing evaluation methods, including analysis, simulation, emulation, testbed, and production deployment, also limit direct numerical comparison.

SmartTE tests a narrower question than these systems: with a common two-path topology and periodic maximum-utilization score, what changes when a lower-score policy is augmented with hysteresis and minimum dwell? The principal gap is not absence of centralized TE or hysteresis in prior work; it is the local measurement of this particular feedback loop across successful route changes and application outcomes under a declared emulated workload. The M5 campaign provides implementation-specific observations only. Broader traffic matrices, more paths and flows, update-consistency mechanisms, hardware tests, and a pre-specified statistical analysis remain future work.

\subsection{Conclusion and future scope}
The literature shows that centralized TE must balance broader demand awareness against telemetry cost, stale state, route stability, and update safety. SmartTE tests only a small emulated feedback-control case and makes no production or universal-superiority claim. Its 19 sources include 13 peer-reviewed papers and two original standards (15 eligible sources), two informational RFCs, one workshop paper not counted as peer reviewed, and one non-peer-reviewed editorial.
